\section{The Communion (nn.~1113--1131)}

With \latincue{Per omnia saecula saeculorum} at the end of the doxology the
priest's voice returns, and the missal returns with it to the \work{Ordo
Missae}. Nineteen numbers carry the rite from the Lord's Prayer to the
ablutions. They are not a single movement but four: a prayer said aloud, a rite
of breaking and peace, a set of private preparations, and the reception itself.
The 1962 book keeps those four visibly distinct, and the distinction is the
best guide to what it thinks Communion is.

\subsection*{The Lord's Prayer and its introduction (nn.~1113--1116)}

\begin{appointed}{1113--1114 $\cdot$ \latincue{Praeceptis salutaribus} and \latincue{Pater noster}}
\upshape\textbf{\vsign.}\ \itshape Per omnia saecula saeculorum.
\upshape\textbf{R.}\ \itshape Amen.\par
\vspace{0.35em}
\upshape(Iungit manus.)\itshape\ Oremus: Praeceptis salutaribus moniti, et
divina institutione formati, audemus dicere:\par
\vspace{0.35em}
\upshape(Extendit manus.)\itshape\ Pater noster, qui es in caelis:
Sanctificetur nomen tuum: Adveniat regnum tuum: Fiat voluntas tua, sicut in
caelo, et in terra. Panem nostrum cotidianum da nobis hodie: Et dimitte nobis
debita nostra, sicut et nos dimittimus debitoribus nostris. Et ne nos inducas
in tentationem. \upshape\textbf{R.}\ \itshape Sed libera nos a malo.
\upshape(Sacerdos secrete dicit:)\itshape\ Amen.
\end{appointed}

\begin{englishwitness}{xxxix}
``Being instructed by thy saving precepts, and following thy divine directions,
we presume to say: Our Father, who art in heaven, hallowed be thy name: thy
kingdom come: thy will be done on earth, as it is in heaven: give us this day
our daily bread; and forgive us our trespasses, as we forgive them that
trespass against us. And lead us not into temptation: \textbf{R.} But deliver
us from evil. \textbf{P.} Amen.''
\end{englishwitness}

Three facts about how the 1962 book prints this text govern everything that can
be said about it.

The first is that the Lord's Prayer is \emph{sung or said aloud, and by the
priest alone}. The general rubrics list \latincue{oratio dominica cum sua
praefatione} among the things said in a clear voice at a read Mass (n.~511 i),
and the \work{Ordo Missae} prints the ministers' part as a single line,
\latincue{Sed libera nos a malo}. The people do not say the Our Father at Mass
in this rite. They answer its last petition, and the priest's \latincue{Amen} to
that answer is said secretly. Whatever one thinks of the arrangement, it is the
book's, and it is stated twice --- in the rubrics and in the layout.

The second is the \term{introduction}, which the missal treats as part of the
prayer and not as a rubrical cue: \latincue{Praeceptis salutaribus moniti, et
divina institutione formati, audemus dicere}. Warned by saving precepts and
formed by divine instruction, we \emph{dare} to say. \latincue{Audemus} is a
strong verb and it is the point of the sentence. The Church does not say the Our
Father at this moment because it is a good prayer; it says it because it was
commanded to, and it says so before saying it. The 1861 English caught the
structure but softened the verb to ``we presume to say'', which is the
translator's difficulty and not the Latin's.

The third is a word. The missal prints \latincue{Panem nostrum
\emph{cotidianum}}; the 1861 hand missal prints \latincue{quotidianum}. That is
the same orthographic restoration met at \latincue{exspecto} in the Creed. But
underneath it lies a real textual choice that is not orthographic at all. The
Vulgate has two forms of the petition: Matthew 6, 11 reads \latincue{panem
nostrum supersubstantialem}, which the Douay-Rheims renders ``Give us this day
our supersubstantial bread'', while Luke 11, 3 reads \latincue{panem nostrum
cotidianum}, ``Give us this day our daily bread''. The Roman Mass says the
Lucan form. It has never used the Matthaean one, and the older commentators who
wanted the eucharistic reading had to argue for it against the text the rite
actually says.

\begin{rubricnote}{Two tones, and where they are used}
The prayer is printed twice with music. The first setting (nn.~1113--1114) is
the solemn tone. The second (nn.~1115--1116) carries the rubric that it is used
on ferias, on vigils of the second and third class, at fourth-class votive
Masses, and at Requiems --- the same austere class of day at which the
\latincue{Gloria Patri} disappears and the \latincue{Ite, missa est} takes its
plainest melody. The chant, in this book, is a calendar in sound.
\end{rubricnote}

Fortescue records the tradition, which he takes from Gregory the Great's own
letter, that Gregory moved the Lord's Prayer to its present place immediately
after the Canon and before the fraction; it is one of three changes he names as
Gregory's, and it is the reason the same page can say both that Gregory ``was
the last to touch the essential part of the Mass'' and that he touched
it.\footnote{Fortescue, \work{The Mass}, p.~172, read at the page image.}

\subsection*{The embolism, the fraction, and the peace (nn.~1117--1119)}

\begin{appointed}{1117 $\cdot$ \latincue{Libera nos}, with the fraction}
\upshape(Deinde manu dextera accipit inter indicem et medium digitos patenam,
quam tenens super altare erectam, dicit secrete:)\itshape\par
Libera nos, quaesumus Domine, ab omnibus malis, praeteritis, praesentibus et
futuris: et intercedente beata et gloriosa semper Virgine Dei Genetrice Maria,
cum beatis Apostolis tuis Petro et Paulo, atque Andrea, et omnibus Sanctis,
\upshape(signat se cum patena a fronte ad pectus,)\itshape\ da propitius pacem
in diebus nostris: \upshape(patenam osculatur,)\itshape\ ut, ope misericordiae
tuae adiuti, et a peccato simus semper liberi et ab omni perturbatione
securi.\par
\vspace{0.35em}
\upshape(\dots\ accipit hostiam, et eam super calicem tenens utraque manu,
frangit per medium, dicens:)\itshape\ Per eundem Dominum nostrum Iesum
Christum, Filium tuum. \upshape(\dots\ frangit particulam, dicens:)\itshape\
Qui tecum vivit et regnat in unitate Spiritus Sancti Deus.\par
\vspace{0.35em}
\upshape\textbf{\vsign.}\ \itshape Per omnia saecula saeculorum.
\upshape\textbf{R.}\ \itshape Amen.
\end{appointed}

\begin{englishwitness}{xxxix--xl}
``Deliver us, we beseech thee, O Lord, from all evils, past, present, and to
come; and by the intercession of the blessed and ever glorious Virgin Mary
Mother of God, and of the holy apostles Peter and Paul, and of Andrew, and of
all the saints, mercifully grant peace in our days: that thro' the assistance
of thy mercy, we may be always free from sin, and secure from all disturbance.
Through the same Jesus Christ, thy Son our Lord, who with thee and the Holy
Ghost liveth and reigneth, God. \textbf{P.} World without end. \textbf{R.}
Amen.''
\end{englishwitness}

\latincue{Libera nos} is an embolism: an insertion that takes the last petition
of the Lord's Prayer and unfolds it. \latincue{Sed libera nos a malo} has just
been sung by the ministers; the priest begins \latincue{Libera nos, quaesumus
Domine, ab omnibus malis} and specifies the evils as past, present and to come.
The prayer is therefore not a new subject but a gloss, and this is why it is
said silently: it is the priest continuing, under his breath, a sentence the
church has just heard aloud.

Two features repay attention. The first is the list of intercessors, which is
short and peculiar: the Virgin, Peter and Paul, and Andrew. Andrew's presence
beside the two Roman apostles, with no other name, is unexplained in the book.
The second is what is asked --- \latincue{da propitius pacem in diebus nostris},
and then two negative goods, freedom from sin and security from all
disturbance. It is the plainest petition in the Mass and the only one that asks
for quiet.

The fraction happens inside this prayer, at its conclusion. The priest breaks
the Host in half at \latincue{Per eundem Dominum nostrum Iesum Christum, Filium
tuum} and breaks a small particle from the left-hand half at \latincue{Qui
tecum vivit et regnat}. The concluding formula of the embolism is thus divided
between two acts of breaking, and it is the only place in the Roman Mass where
a conclusion is used as a rubric.

\begin{appointed}{1118--1119 $\cdot$ \latincue{Pax Domini} and \latincue{Haec commixtio}}
\upshape(Cum ipsa particula signat ter super calicem, dicens:)\itshape\par
Pax \misscross\ Domini sit \misscross\ semper vobis\misscross cum.
\upshape\textbf{R.}\ \itshape Et cum spiritu tuo.\par
\vspace{0.35em}
\upshape(Particulam ipsam immittit in calicem, dicens secrete:)\itshape\par
Haec commixtio, et consecratio Corporis et Sanguinis Domini nostri Iesu
Christi, fiat accipientibus nobis in vitam aeternam. Amen.
\end{appointed}

\begin{englishwitness}{xl}
``May the peace of the Lord be always with you. \textbf{R.} And with thy
spirit.'' And, at the commingling: ``May this mixture, and consecration of the
body and blood of our Lord Jesus Christ, be to us that receive it effectual to
eternal life. Amen.''
\end{englishwitness}

The greeting is not \latincue{Dominus vobiscum} but \latincue{Pax Domini sit
semper vobiscum}, and it is used once in the Mass, here. It is made with the
particle of the Host held over the chalice, and three crosses are signed with
it. The particle is then dropped into the chalice while the priest says
\latincue{Haec commixtio}.

The commingling is one of the harder things in the Roman rite to state
accurately, and the safest course is to say what the book says and no more. The
formula calls what is happening a \latincue{commixtio} and a
\latincue{consecratio} --- a mixing and a hallowing --- and asks that it be
\latincue{accipientibus nobis in vitam aeternam}, unto eternal life for us who
receive. It offers no account of why the two species are recombined, and the
1962 books nowhere supply one. The medieval explanations are many and mutually
inconsistent; none of them is in this missal, and a study of this missal should
not import one.

\subsection*{\latincue{Agnus Dei} (n.~1120)}

\begin{appointed}{1120 $\cdot$ \latincue{Agnus Dei}}
Agnus Dei, qui tollis peccata mundi: miserere nobis.\par
Agnus Dei, qui tollis peccata mundi: miserere nobis.\par
Agnus Dei, qui tollis peccata mundi: dona nobis pacem.
\end{appointed}

\begin{rubricnote}{Posture, and the Requiem form}
The priest bows toward the Sacrament with joined hands and strikes his breast
three times, saying the invocations in an intelligible voice. In Requiem Masses
\latincue{miserere nobis} is not said; in its place \latincue{dona eis requiem},
and in the third invocation \latincue{sempiternam} is added, so that the last
line becomes \latincue{dona eis requiem sempiternam}. The 1861 hand missal
prints exactly this variant in English --- ``Give them rest'', and ``Give them
eternal rest'' --- which shows the rule is older than the code of 1960.
\end{rubricnote}

The words come from the Baptist's greeting at the Jordan, and not verbatim: the
Vulgate at John 1, 29 has \latincue{Ecce agnus Dei, ecce qui tollit peccatum
mundi}, in the third person and with sin in the singular, which the
Douay-Rheims renders ``Behold the Lamb of God. Behold him who taketh away the
sin of the world.'' The Mass turns the sentence into an address and the noun
into a plural. That is a small change and a real one: the liturgical form asks
mercy of the Lamb rather than pointing him out, and it names sins rather than
sin.

The address has already been used once in the same Mass, in the
\latincue{Gloria} at \latincue{Qui tollis peccata mundi, miserere nobis}. Here
it is used three times and the third ending changes: \latincue{dona nobis
pacem}. The change is the reason the formula stands at this point of the rite
rather than at any other, for the kiss of peace at solemn Mass follows
immediately.

\begin{witnesscard}{Historical note: what the \latincue{Agnus Dei} was for}
Fortescue's account is that the chant ``was added to fill up the time of the
fraction'': the \work{Liber Pontificalis} says Sergius I (687--701) ordered it
to be sung by clergy and people at the breaking of the Lord's body; it was at
first sung once, twice by the eleventh century, and John Beleth in the twelfth
describes the present threefold form ending \latincue{dona nobis pacem}, while
Innocent III notes churches, the Lateran among them, which kept
\latincue{miserere nobis} three times.\footnote{Fortescue, \work{The Mass},
p.~387, read at the page image.} If that account is right, then the 1962
arrangement --- the chant sung while the priest has already finished breaking
--- is a survival of a longer fraction, and the Requiem form, which keeps three
identical endings, is the older shape of the piece.
\end{witnesscard}

\subsection*{The three prayers before Communion (nn.~1121--1124)}

\begin{appointed}{1121--1122 $\cdot$ \latincue{Domine Iesu Christe, qui dixisti} and the peace}
Domine Iesu Christe, qui dixisti Apostolis tuis: Pacem relinquo vobis, pacem
meam do vobis: ne respicias peccata mea, sed fidem Ecclesiae tuae; eamque
secundum voluntatem tuam pacificare et coadunare digneris: Qui vivis et regnas
Deus per omnia saecula saeculorum. Amen.\par
\vspace{0.35em}
\upshape(Si danda est pax, osculatur altare et, dans pacem, dicit:)\itshape\
Pax tecum. \upshape\textbf{R.}\ \itshape Et cum spiritu tuo.
\end{appointed}

\begin{rubricnote}{The Requiem exception}
The missal prints, immediately after: \latincue{In Missis defunctorum non datur
pax, neque dicitur praecedens oratio.} At a Requiem the peace is not given and
this prayer is not said. The other two preparatory prayers are.
\end{rubricnote}

The first prayer is about the Church and not about the priest, which is why it
is the one omitted when there is no kiss of peace to give. It quotes John 14,
27 --- \latincue{Pacem relinquo vobis, pacem meam do vobis} --- and then makes
the request that gives the sentence its force: \latincue{ne respicias peccata
mea, sed fidem Ecclesiae tuae}. Do not look at my sins, but at your Church's
faith. The priest asks God to look away from the very man who is speaking.

\begin{appointed}{1123--1124 $\cdot$ \latincue{Domine Iesu Christe, Fili Dei vivi} and \latincue{Perceptio}}
Domine Iesu Christe, Fili Dei vivi, qui ex voluntate Patris, cooperante Spiritu
Sancto, per mortem tuam mundum vivificasti: libera me per hoc sacrosanctum
Corpus et Sanguinem tuum ab omnibus iniquitatibus meis et universis malis: et
fac me tuis semper inhaerere mandatis, et a te numquam separari permittas: Qui
cum eodem Deo Patre et Spiritu Sancto vivis et regnas Deus in saecula
saeculorum. Amen.\par
\vspace{0.35em}
Perceptio Corporis tui, Domine Iesu Christe, quod ego indignus sumere praesumo,
non mihi proveniat in iudicium et condemnationem: sed pro tua pietate prosit
mihi ad tutamentum mentis et corporis, et ad medelam percipiendam: Qui vivis et
regnas cum Deo Patre in unitate Spiritus Sancti Deus, per omnia saecula
saeculorum. Amen.
\end{appointed}

\begin{englishwitness}{xli}
``Lord Jesus Christ, Son of the living God, who, according to the will of thy
Father, hast by thy death, thro' the co-operation of the Holy Ghost, given life
to the world, deliver me by this thy most sacred body and blood from all my
iniquities, and from all evils; and make me always adhere to thy commandments,
and never suffer me to be separated from thee \dots'' And: ``Let not the
participation of thy body, O Lord Jesus Christ, which I, though unworthy,
presume to receive, turn to my judgment and condemnation, but thro' thy mercy,
may it be a safeguard and remedy, both to soul and body \dots''
\end{englishwitness}

These two are personal in a way nothing else in the Mass is. Every verb is
singular and first person --- \latincue{libera me}, \latincue{fac me},
\latincue{prosit mihi} --- and their subject is the danger of the act about to
be performed. \latincue{Perceptio} states it in the harshest available terms:
that receiving may turn out \latincue{in iudicium et condemnationem}, into
judgment and condemnation. That is 1 Corinthians 11, 29 turned into a prayer
against itself. The Roman rite lets the priest name the possibility of eating
to his own harm at the moment before he eats, and asks instead for
\latincue{tutamentum mentis et corporis}, a safeguard of mind and body, and for
\latincue{medela}, a cure.

The trinitarian clauses at the end of the two prayers are worth a glance for
what they show about how carefully this Latin was assembled. The first ends
\latincue{Qui cum eodem Deo Patre et Spiritu Sancto vivis et regnas} ---
\latincue{eodem} because the Father has been named inside the prayer --- and the
second ends \latincue{Qui vivis et regnas cum Deo Patre in unitate Spiritus
Sancti}, without \latincue{eodem}, because he has not. The same scruple runs
through the Canon.

\subsection*{The reception (nn.~1125--1129)}

\begin{appointed}{1125--1127 $\cdot$ \latincue{Panem caelestem}, \latincue{Domine non sum dignus}, and the Body}
Panem caelestem accipiam, et nomen Domini invocabo.\par
\vspace{0.35em}
\upshape(\dots\ dextera tribus vicibus percutiens pectus, elata aliquantulum
voce, ter dicit devote et humiliter:)\itshape\par
Domine, non sum dignus, \upshape(et secrete prosequitur:)\itshape\ ut intres
sub tectum meum: sed tantum dic verbo, et sanabitur anima mea.\par
\vspace{0.35em}
\upshape(Postea, dextera se signans cum hostia super patenam, dicit:)\itshape\par
Corpus Domini nostri Iesu Christi custodiat animam meam in vitam aeternam.
Amen.
\end{appointed}

\begin{englishwitness}{xlii}
``Lord, I am not worthy that thou shouldst enter under my roof; say but the
word, and my soul shall be healed.'' And: ``May the body of our Lord Jesus
Christ preserve my soul to life everlasting. Amen.''
\end{englishwitness}

\latincue{Domine, non sum dignus} is the centurion's answer at Matthew 8, 8, and
the alteration made to it is exact and small. The Douay-Rheims has the centurion
say ``Lord, I am not worthy that thou shouldst enter under my roof; but only
say the word, and my servant shall be healed.'' The Mass changes
\latincue{puer meus}, my servant, to \latincue{anima mea}, my soul, and changes
nothing else. The man who asked for a cure at a distance becomes the man asking
for one within.

The rubric divides the sentence between two voices, and the division is the
same one used at \latincue{Nobis quoque peccatoribus} in the Canon: four words
raised, the rest secret. In a rite where almost everything near the altar is
inaudible, the two things the people are given to hear are an admission of
sinfulness and an admission of unworthiness.

\begin{appointed}{1128--1129 $\cdot$ \latincue{Quid retribuam} and the Blood}
Quid retribuam Domino pro omnibus quae retribuit mihi? Calicem salutaris
accipiam, et nomen Domini invocabo. Laudans invocabo Dominum, et ab inimicis
meis salvus ero.\par
\vspace{0.35em}
Sanguis Domini nostri Iesu Christi custodiat animam meam in vitam aeternam.
Amen.
\end{appointed}

The missal prints no scriptural reference here, though it prints one at the
\latincue{Lavabo} and at \latincue{Dirigatur}. The reason is probably that the
text is not one passage. \latincue{Quid retribuam Domino \dots\ et nomen Domini
invocabo} is Psalm 115, 12--13 in the Vulgate numbering --- ``What shall I
render to the Lord, for all the things that he hath rendered to me? I will take
the chalice of salvation; and I will call upon the name of the Lord'' --- and
\latincue{Laudans invocabo Dominum, et ab inimicis meis salvus ero} is Psalm
17, 4, ``Praising, I will call upon the Lord: and I shall be saved from my
enemies.'' The rite has joined a verse about drinking a cup to a verse about
being rescued, and the joint is invisible in the Latin.

The two reception formulas are identical but for one word, \latincue{Corpus}
and \latincue{Sanguis}, and both ask the same thing: that what is received
\latincue{custodiat animam meam in vitam aeternam}, keep my soul unto life
everlasting. \latincue{Custodire} is to guard, not to bring; the petition is for
preservation, and it matches the two preparatory prayers, which asked for a
safeguard rather than for an increase.

\subsection*{The Communion of the people, which the \work{Ordo Missae} does not print}

This is the most consequential structural fact in the whole Communion rite, and
it is easy to miss because it consists of an absence.

\begin{rubricnote}{\work{Rubricae generales Missalis romani} nn.~502--503}
n.~502: the proper time for giving Holy Communion to the faithful is within
Mass, after the Communion of the celebrant, who is himself to distribute it to
those who ask, unless the number of communicants makes help from another priest
appropriate; it is altogether unfitting that Communion be distributed at the
same altar at which Mass is actually being celebrated by another priest outside
the proper time; and for a reasonable cause Communion may also be given
immediately before or after Mass, or outside Mass, in which cases the form
prescribed in the \work{Rituale Romanum}, title V, chapter II, nn.~1--10, is
used. n.~503: whenever Holy Communion is distributed within Mass, the celebrant,
having consumed the Precious Blood, proceeds immediately to the distribution,
\emph{the confession and absolution being omitted}, but \latincue{Ecce Agnus
Dei} and the threefold \latincue{Domine, non sum dignus} being said.
\end{rubricnote}

The \work{Ordo Missae} prints, at the end of n.~1129, one clause about the
people: \latincue{Quo sumpto, si qui sunt communicandi, eos communicet,
antequam se purificet} --- when the Precious Blood has been consumed, if there
are any to be communicated, let him communicate them before he purifies himself.
That is the whole of it. No formula is printed. \latincue{Ecce Agnus Dei},
\latincue{Domine, non sum dignus} said three times by or for the communicants,
and \latincue{Corpus Domini nostri Iesu Christi custodiat animam tuam in vitam
aeternam} are prescribed by the general rubrics and printed in another book.

Two things follow and both should be said plainly. First, the Communion of the
faithful in this Missal is legislated but not typeset: a reader who takes the
\work{Ordo Missae} as the complete order of the rite will not find it. Second,
n.~503 is a piece of recent legislation and not an ancient custom. The
\latincue{Confiteor} before the people's Communion --- the servers' confession
repeated at the rail --- is precisely what n.~503 suppresses, and its
suppression is one of the visible marks of the 1960 code in a rite otherwise
unchanged since 1570. What replaced it is nothing; the priest turns from the
altar and says \latincue{Ecce Agnus Dei}.

\subsection*{The ablutions (nn.~1130--1131)}

\begin{appointed}{1130--1131 $\cdot$ \latincue{Quod ore sumpsimus} and \latincue{Corpus tuum}}
Quod ore sumpsimus, Domine, pura mente capiamus: et de munere temporali fiat
nobis remedium sempiternum.\par
\vspace{0.35em}
Corpus tuum, Domine, quod sumpsi, et Sanguis, quem potavi, adhaereat visceribus
meis: et praesta; ut in me non remaneat scelerum macula, quem pura et sancta
refecerunt sacramenta: Qui vivis et regnas in saecula saeculorum. Amen.
\end{appointed}

\begin{englishwitness}{xliii}
``May thy body, O Lord, which I have received, and thy blood which I have
drank, cleave to my bowels; and grant that no stain of sin may remain in me,
who have been fed with this pure and holy sacrament. Who livest and reignest,
\&c.''
\end{englishwitness}

The first prayer is in the plural and the second in the singular, and the
difference tracks the two ablutions: the first is said while the priest gathers
any fragments and wipes the paten over the chalice, before the wine is poured;
the second while he purifies his fingers. \latincue{Quod ore sumpsimus \dots\
pura mente capiamus} sets \latincue{ore} against \latincue{mente} and
\latincue{sumere} against \latincue{capere}, and asks that what was taken by
mouth be grasped by mind. \latincue{De munere temporali fiat nobis remedium
sempiternum} is the same movement again: a gift given in time, a remedy that
does not end.

\latincue{Corpus tuum \dots\ adhaereat visceribus meis} is the most physical
sentence in the Roman Mass, and the 1861 translator did not flinch at it:
``cleave to my bowels''. \latincue{Viscera} is the inward parts. The prayer asks
that what has been eaten stay, and then names the reason it might not: that a
stain of sins should not remain in a man whom pure and holy sacraments have
fed. The last relative clause, \latincue{quem pura et sancta refecerunt
sacramenta}, is in the accusative and hangs on \latincue{in me}; the Latin is
awkward and the awkwardness is old.

After this the priest wipes the chalice, folds the corporal, covers the chalice
and replaces it on the altar as before, and goes on with the Mass. The
Communion antiphon --- proper, and printed with the day's Mass --- and the
postcommunion follow, in the same number and order as the collects, with the
Prayer over the People added on the ferias of Lent and Passiontide outside the
Sacred Triduum (nn.~504--506).
